\documentclass[11pt]{article}
\usepackage[margin=1in]{geometry}
\usepackage{fontspec,unicode-math,amsmath}
\setmainfont[Path=fonts/,BoldFont=lmroman10-bold.otf,ItalicFont=lmroman10-italic.otf,BoldItalicFont=lmroman10-bolditalic.otf]{lmroman10-regular.otf}
\setmonofont[Path=fonts/,Scale=.86,BoldFont=DejaVuSansMono-Bold.ttf,ItalicFont=DejaVuSansMono-Oblique.ttf,BoldItalicFont=DejaVuSansMono-BoldOblique.ttf]{DejaVuSansMono.ttf}
\setmathfont[Path=fonts/]{latinmodern-math.otf}
\usepackage{booktabs,longtable,array,fvextra,microtype,xcolor,fancyhdr}
\usepackage[colorlinks=true,linkcolor=blue,urlcolor=blue]{hyperref}
\special{pdf:trailerid [<20261006000000000000000000000012><20261006000000000000000000000012>]}
\special{pdf:docinfo << /CreationDate (D:20261006000000Z) /ModDate (D:20261006000000Z) >>}
\hypersetup{pdftitle={Paper XII: Final Lean Guide},pdfauthor={Komorebi P.},pdfsubject={Production-source guide, exact hypotheses, claim crosswalk and replay scope}}
\setlength{\parskip}{5pt}\setlength{\parindent}{0pt}
\setlength{\emergencystretch}{4em}
\pagestyle{fancy}\fancyhf{}\fancyhead[L]{\small Paper XII: Final Lean Guide}\fancyhead[R]{\small Komorebi P.}\fancyfoot[C]{\thepage}
\newcommand{\code}[1]{\nolinkurl{#1}}
\newcommand{\module}[1]{\nolinkurl{PaperXII.#1}}
\Urlmuskip=0mu plus 1mu
\makeatletter
\g@addto@macro\UrlBreaks{\do\a\do\b\do\c\do\d\do\e\do\f\do\g\do\h\do\i\do\j\do\k\do\l\do\m\do\n\do\o\do\p\do\q\do\r\do\s\do\t\do\u\do\v\do\w\do\x\do\y\do\z\do\A\do\B\do\C\do\D\do\E\do\F\do\G\do\H\do\I\do\J\do\K\do\L\do\M\do\N\do\O\do\P\do\Q\do\R\do\S\do\T\do\U\do\V\do\W\do\X\do\Y\do\Z\do\0\do\1\do\2\do\3\do\4\do\5\do\6\do\7\do\8\do\9}
\makeatother
\newcommand{\excerpt}[3]{\par\smallskip\textit{Source: \code{PaperXII/#1.lean}, lines #2--#3.}\par\VerbatimInput[fontsize=\scriptsize,breaklines=true,breakanywhere=true,numbers=left,numbersep=4pt,firstline=#2,lastline=#3,firstnumber=#2,frame=lines]{excerpts/#1.display.txt}\par}
\newcommand{\Supp}{\operatorname{Supp}}\newcommand{\Spec}{\operatorname{Spec}}
\newcommand{\Ann}{\operatorname{Ann}}\newcommand{\Fitt}{\operatorname{Fitt}}
\newcommand{\OO}{\mathcal O}\newcommand{\TT}{\mathbb T}
\newcommand{\wideotimes}{\mathbin{\widehat\otimes}}
\title{Paper XII: Final Lean Guide\\\large Patching Support, Hecke Packets, and Conditional Automorphic Completion}
\author{Komorebi P.}\date{2026-10-06}

% Operation Komorebi: shared publication profile.
\usepackage{komorebi-publication}
\renewcommand{\KomorebiTitle}{Patching Support and Hecke Packets}
\renewcommand{\KomorebiSubtitle}{Conditional automorphic completion}
\renewcommand{\KomorebiSeries}{Core series XII}
\renewcommand{\KomorebiKind}{Lean and evidence guide}
\renewcommand{\KomorebiEdition}{Standardized publication edition / 8 October 2026}
\renewcommand{\KomorebiShortTitle}{Paper XII}
\renewcommand{\KomorebiTitleSize}{31}
\renewcommand{\KomorebiTitleLeading}{36}
\renewcommand{\KomorebiProfile}{classical}
\renewcommand{\KomorebiOrientation}{portrait}
\KomorebiSetup
\begin{document}
\KomorebiCover

\relax
\begin{abstract}
This guide explains the complete production Lean source accompanying Paper XII, with a source catalog, declaration census, import audit, exact source excerpts, and a claim crosswalk. It separates constructed algebraic and topological objects from hypotheses supplied by an arithmetic application. The checked revised contract comprises 49 internal claims mapped to 291 distinct production targets. The original wording remains a separate ledger: 41 complete, seven partial, and one unchecked. The guide preserves all 96 production modules and their inherited verification records. Its new checks are source, provenance, and document checks; no fresh Lean compilation was available in the finalization environment.
\end{abstract}
\tableofcontents
\newpage

\section{What this guide certifies and how to use it}
The reference production tree is the \code{formal_release/lean/} tree from the complete v0.17.1 bundle. Its revised manuscript is the 36-page \code{Paper_XII_Revised_Formalization_Complete_v0.17.0.tex}; the v0.17.1 release reverified its unchanged production sources. This guide is dated 2026-10-06 and describes those exact sources, alongside the final manuscript's explicit editorial clarifications. It does not silently replace any production theorem or historical verification record.

Read Sections 4--7 for the completed local and global constructions, Sections 8--12 for support, packets, occupancy and conditional realization, and Section 14 for replay. Appendices list every source module and every revised internal claim. The machine-readable catalog gives full SHA-256 hashes, source locations, imports, declaration names and line numbers. The full production source is supplied; printing every line in this PDF is unnecessary for reproducibility.

Three distinct kinds of evidence must be kept separate. A Lean kernel check verifies the exact formal proposition under its hypotheses and dependencies. An ordinary mathematical proof explains a proposition in mathematical prose, with its own statement and citation obligations. A finite computation verifies the tested cases, numerical identities or model bounds stated by that computation. Neither a finite numerical certificate nor a source-token scan is a substitute for kernel elaboration of the production proof.

\begin{center}
\begin{tabular}{p{.26\textwidth}p{.64\textwidth}}
\toprule Evidence & Exact scope\\\midrule
Inherited Lean replay & September 12, 2026: all 96 production modules compiled; 1,229 named declarations axiom-audited; three invalid controls rejected. Pinned Mathlib cache used.\\
Finalization source audit & Production hashes, lexical census, direct imports, claim targets, source preservation, and inherited report consistency checked on 2026-10-06.\\
New guide build & Portable PDF build, embedded-font/glyph/reference checks, and independent fresh-directory PDF/text/page comparison.\\
Numerical verification & Separate finite and exact computation records explain each tested input and bound. No universal arithmetic statement follows from those tests.\\
\bottomrule
\end{tabular}
\end{center}

\textbf{No fresh Lean compilation is claimed in this finalization.} The pinned Lean executable, Lake and Elan were unavailable in the finalization environment. Preserved compiler logs and matching source/object hashes are inherited evidence, not a new run of the kernel. The replay commands below enable a reader with the separately supplied pinned runtime to make such a run.

\section{Independently checked counts and their definitions}
The registered manifest contains exactly 96 distinct modules, and its source path set equals the set of all production \code{.lean} files. There are 91 files under \code{PaperXII/} and five under \code{PaperX/}. Directory names alone do not identify originality: \module{ReuseCompactness} and \module{ReuseFiniteSeparation} retain Paper X namespaces and 18 reused theorem declarations. Together with 63 theorem declarations in the five Paper X files, these account for the 81 reused theorems. Paper XII's own public theorem count is 748.

The public named lexical census is 1,229: 829 theorems, 319 definitions, 47 abbreviations, 17 structures and 17 named instances. One private theorem helper, \code{PaperXII.DeterminantalIdeals.option_factor}, is excluded from the public count. The 829 theorem declarations divide into 748 own and 81 reused; they are supporting lemmas and interfaces as well as theorems implementing numbered manuscript claims. This is not a historical novelty count or a twelve-paper total.

The independent scanner blanks nested block comments, line comments and strings while preserving line numbers. It recognizes anchored named declarations in each module's unique namespace and resolves their qualified names. Its declaration set agrees exactly with the inherited axiom target set. It finds no production \code{sorry}, \code{admit}, \code{axiom}, \code{unsafe}, or \code{native_decide} token outside comments/strings. These are lexical findings; elaboration and transitive axioms remain the role of the compiler records.

The 49 revised claim records comprise 27 theorems, 14 corollaries, seven propositions and one lemma. Three worked examples, definitions, remarks and eight external standard facts are outside that internal census. All 291 distinct mapped targets occur in the 1,229-name inherited audit. A target can support several claims; a claim can require several targets. Target count is therefore neither theorem count nor claim count.

The original ledger has 41 complete, seven partial and one unchecked claim. Its eight inherited amendments are 6.12, 9.4, 9.5, 9.6, 10.5, 12.2, 12.4 and 12.6. The revised coverage is 49 complete with zero partial and zero unchecked. Both ledgers are supplied without relabeling the original statements as proved. The former 47/49 textual retraction is superseded by the recovered full production archive and successful September replay; it is not an outstanding formal obligation.

\section{Source organization, dependency order and hypothesis notation}
The source audit finds 113 direct internal import edges. Every internal import names a registered module and precedes its consumer in \code{checks/modules.json}. The external direct imports are Mathlib modules and \code{Std}; the complete list is in \code{data/census.json}. This checks source-level dependency order, not the entire transitive Mathlib import graph. The production tree reuses Paper X, not Paper XI.

The original 71-module block, ten Fitting modules and 15 completion modules make up the 96-module cumulative release. The 15-module completion block contains 146 public theorem declarations and 225 public named declarations. These are preserved existing work, not newly proved results in this guide. Source preservation and the declaration census are computed against the actual files.

In Lean, square-bracket arguments such as \code{[CommRing R]} and \code{[CompleteSpace P]} are typeclass hypotheses. A proposition-valued ordinary argument such as \code{(hB : IsAdic mB)} is also a hypothesis. A \code{variable} command may contribute ambient arguments to later declarations; \code{include} and \code{omit} control that contribution. Reading only the line beginning \code{theorem} can therefore hide essential assumptions. The full exact excerpt modules preserve these commands, imports, namespace and section boundaries.

The source excerpts delivered as \code{excerpts/*.lean} are byte-identical full modules, including their ambient hypotheses. PDF excerpts select line ranges and retain original line numbers. To avoid missing glyphs in portable code typography, the display copies replace eleven symbols by readable forms, for example \code{iSup} for the indexed supremum and \code{/} for Lean's ideal-quotient slash. The complete replacement map is in \code{excerpts/INDEX.json}. This is a documented display convention only; production source and exact excerpts retain every original Unicode character.

\section{Claim 3.1: construct the tensor, then complete it}
For a finite index type \(\iota\), commutative coefficient ring \(O\), and commutative \(O\)-algebras \(R_i\), the raw tensor is Mathlib's actual \code{PiTensorProduct O R}. It is not a record accepting an already supplied tensor equivalence. The factor inclusion is \code{PiTensorProduct.singleAlgHom}; the family lift is constructed from a multilinear product and the tensor's universal map.

Given ideals \(m_O\subseteq O\), \(m_i\subseteq R_i\), and condition ideals \(I_i\), define
\[
S=\bigotimes_{i\in\iota,O}R_i,\qquad
\mathfrak a=m_OS+\sum_i m_iS,\qquad L=\sum_i I_iS.
\]
In Lean the finite sums are indexed suprema of mapped ideals. The scalar term is present even if \(\iota\) is empty. Removing it would change the empty-family topology to the wrong normalization.
\excerpt{FiniteTensorQuotient}{11}{26}
\excerpt{FiniteTensorQuotient}{60}{75}

The ordinary quotient map is built by applying the ideal quotient in each factor. Its kernel contains \(L\); quotient lifts construct both directions between \(S/L\) and \(\bigotimes_i(R_i/I_i)\). Tensor extensionality and quotient extensionality prove the inverse identities. This produces an ordinary algebra equivalence, with a proved exact image for the defining ideal. Ordinary tensor right exactness alone would not establish the completed theorem.

\code{AdicModel.Carrier R I} is a type synonym that installs the specified ideal-adic topology on \(R\). Its \code{Completed} object is the actual Hausdorff uniform completion of that carrier. The forgetful raw ring equivalence does not assert continuity into an arbitrarily chosen topology. This prevents the topology from being replaced silently.
\excerpt{AdicModel}{12}{42}

Thus \code{CompletedTensor R mO m} is the Hausdorff completion of \(S\) in the literal \(\mathfrak a\)-adic additive uniformity. Its \(O\)-algebra structure comes from the dense embedding of \(S\), and its completed factor maps compose that embedding with the raw factor inclusions. The condition ideal is the closure of the algebraically extended \(L\).
\excerpt{CompletedFiniteTensor}{16}{47}

\section{Completion and quotient: actual maps, completeness and inverses}
The reusable comparison is
\[
\widehat R\big/\overline{I\widehat R}\ \cong\ \widehat{R/I}.
\]
Its ambient assumptions are a commutative ring, a uniform space with uniform additive group, and a topological ring. Countable generation of the additive uniformity is used to prove completeness of the quotient of the completion. There is no assumption supplying completeness of that quotient, an inverse comparison map, or comparison surjectivity. The closure makes the quotient Hausdorff.
\excerpt{CompletionQuotient}{18}{61}

The forward map extends the actual quotient map. Its continuity and the target's separation prove that it kills the closed extended ideal. The backward map first descends the dense raw representatives, then extends from \(R/I\) to its completion. The source quotient has its actual quotient topology and the associated additive uniformity. \module{AdicQuotientTopology} separately proves that this quotient topology is the topology of the image defining ideal.

The inverse laws are proved on dense representatives and extended by continuity. The second law uses quotient surjectivity after completion induction. The ring equivalence and its homeomorphism are then constructed from these proved maps.
\excerpt{CompletionQuotient}{121}{180}

Applying this to the raw finite tensor quotient gives the displayed \(O\)-algebra equivalence of Claim 3.1. Its forward direction in the manuscript is the inverse of the construction direction. Both maps are continuous, and every factor identity is proved with the actual inclusions.
\excerpt{CompletedFiniteTensor}{91}{133}

The source theorem \code{closedConditions_eq} identifies the ideal with the literal closure of the ideal generated by the completed factor images. The supplementary consumer \code{exact_completed_kernel} additionally asserts exact equality of the canonical completed-map kernel and the stated closed ideal. That consumer is separate from the production theorem census.

The general comparison requires neither locality nor Noetherianity of the raw tensor ring. \module{AdicCompletionComparison} independently identifies the uniform completion with Mathlib's algebraic inverse-limit adic completion. Finite generation controls ideal powers and identifies the completed topology with the topology of its extended defining ideal. These steps support the subsequent local-ring specialization.

\section{Locality, universal maps and the edge cases}
For the manuscript specialization, supply a field \(k\), an \(O\)-algebra structure on \(k\), residue algebra maps \(R_i\to k\), and surjectivity of \(O\to k\). The defining ideals are their kernels. \module{FiniteTensorLocality} constructs \(S/\mathfrak a\cong k\), proves \(\mathfrak a\) maximal, and derives finite generation from Noetherianity of \(O\) and each \(R_i\). It then proves locality and the maximal-ideal-adic topology of the actual completion. These residue maps are explicit mathematical data, not an implicit assertion that all residue fields somehow agree.
\excerpt{FiniteTensorLocality}{100}{138}

The completed universal map also has explicit topological control. Its target \(P\) is a complete Hausdorff uniform topological ring with \(O\)-algebra structure, \(m_P\)-adic topology, and factor maps carrying both the scalar and factor defining ideals into \(m_P\). The raw family map is proved continuous, extended to the actual completion, and shown unique by density and finite tensor extensionality. The final existence-and-uniqueness theorem is an output of that construction.
\excerpt{CompletedTensorUniversal}{15}{23}
\excerpt{CompletedTensorUniversal}{83}{99}

The empty family retains \(m_O\)'s adic topology. Zero conditions give the identity comparison. A unit condition produces the zero ring and still satisfies the ring comparison. Locality of a quotient is claimed only with an explicit nontriviality premise; the zero ring is not assigned the manuscript's nonzero local ring and residue field. Source boundary checks cover these cases separately.

\section{Claim 3.3: completed global base change and its points}
Let \(A\) be the constructed completed local tensor, let \(J\) be its closed condition ideal, and let \(C\) be the independently completed tensor of \(R_i/I_i\). The \(A\)-algebra structure on \(C\) is built from the proved Claim 3.1 equivalence \(A/J\cong C\). \code{quotientIdentification_ideal} proves that the transported defining ideal is the one defining the actual condition tensor topology.

Given the supplied global ring \(B\) and its structural \(A\)-algebra map, the global condition ring is the actual completion of \(B\otimes_A C\). The transport code first constructs the tensor with \(A/J\), then transports along the proved quotient identification. It is not an assumed pushout with a supplied point-factorization field. For complete \(B\), the actual ring is bicontinuously identified with \(B/\overline{JB}\), and its canonical global map is surjective.

The reusable ring-point theorem assumes finite generation of the scalar and factor defining ideals, \(B\)'s declared adic topology, and an ideal bound for the structural map. Its target is a complete Hausdorff uniform topological ring. It proves that a unique continuous point factor exists precisely when each original local restriction kills \(I_i\).
\excerpt{CompletedLocalConditions}{17}{38}
\excerpt{CompletedLocalConditions}{66}{100}

The manuscript capstone uses Noetherianity and common residue maps to supply finite generation. It requires \(B\) local and the structural algebra map local; the resulting maximal ideal bound is proved internally. Locality is an essential premise here. A bare continuous ring map is not the full printed capstone's input. The final traditional manuscript exposes these inherited conventions directly in Claim 3.3 and spells out the point category rather than relying on conventions elsewhere.
\excerpt{CompletedPaperCapstone}{17}{69}

For continuous local \(O\)-algebra points, the production point theorem additionally has \(B\) complete, actual \(O\)-algebra structures on \(B\) and \(P\), and an original point which is an \(O\)-algebra map reflecting units. The completed global coefficient action is the actual composite through \(B\). Surjectivity of the global map proves that the unique ring factor reflects units; equality on coefficients proves it is an \(O\)-algebra map. Naturality follows from equality after the surjective global map.
\excerpt{CompletedPointCategory}{49}{92}

The ring-point result, local algebra-point result and Noetherian/local specialization are distinct declarations and must be read together. A local deformation test point is distinct from a characteristic-zero field specialization, where a coefficient uniformizer can be a unit; continuity of the specialized representation is separately required by realization. The supplementary \code{claim33} consumes the local point statements using manuscript-level hypotheses, without an extra finite-generation or ideal-bound hypothesis. It does not prove arithmetic representability of a deformation functor. Representability, universal arithmetic deformation representations, and existence of the local quotient rings remain supplied inputs.

\section{Ordinary support, patching and intrinsic Fitting ideals}
For a finite module, production support is Mathlib's actual localization support on the prime spectrum. Lean proves
\[
\Supp_R M=V(\Ann_R M).
\]
Under the Noetherian hypotheses, near faithfulness is equivalent to full support. Quotient support is obtained on corresponding primes using actual quotient modules and localization/Nakayama arguments. Ordinary patching then uses supplied coefficient and module control isomorphisms to descend full support from the patched finite module; finite generation and finite-level full support are outputs.

An arithmetic patching theorem supplying the patched ring, module, control maps and full patched support remains external. The finite-level \(R\)-action on the same carrier must be restriction of its Hecke action through the supplied comparison, as expressed by the ambient scalar-tower hypotheses; derived carriers require the corresponding transport of the actual complex. Two unrelated actions on a carrier do not supply this compatibility. Freeness over an auxiliary algebra is insufficient for full deformation support; the power-series counterexample demonstrates the missing implication. A ring isomorphism \(R\cong\TT\) transports existing support. A surjection alone can erase a component, as in the actual node \(k[[x,y]]/(xy)\) followed by quotienting by \((x)\).

The ten-module Fitting block constructs the determinantal ideals of a finite presentation and proves their invariance under changes of presentation, generator and relation counts. Its intrinsic definition chooses a presentation using classical choice and then proves equality with the minors of every actual finite presentation. No presentation-independence theorem is supplied as a hypothesis.
\excerpt{IntrinsicFitting}{8}{39}

The arbitrary base-change theorem transports each minor under the ring map. The support criterion and all higher strata are stated for finitely presented modules; Noetherianity plus finite generation in the manuscript supplies that premise. For a presentation with \(n\) generators,
\[
\dim_{\kappa(\mathfrak p)}(M\otimes\kappa(\mathfrak p))
 =n-\operatorname{rank} A(\mathfrak p).
\]
For \(1\le r\le n\), the inequality \(\dim\ge r\) is detected by minors of size \(n-r+1\). If \(r>n\), the fiber locus is empty and \(\Fitt_{r-1}(M)=R\). The final ordinary proof explicitly separates this boundary; the formal block already includes upper-index normalization and empty/zero-minor checks. Fitting strata count minimal fiber generators, not automorphic multiplicities.

\section{Derived support uses actual perfect models}
\module{PerfectDerivedSupport} does not encode a complex by a bare list of supports. Its \code{PerfectComplex} contains an actual Mathlib derived-category object, a bounded finite-projective cochain model, and an actual derived isomorphism from that model. Base change is computed degree by degree on the projective model. The code proves compatibility with actual derived isomorphisms, so support does not depend on the chosen model.
\excerpt{PerfectDerivedSupport}{23}{53}

Derived support is the set of primes with nonzero residue-field derived fiber. Scalar extension has spectral inverse-image support; the quotient identity is expressed as an image in the original spectrum. It includes the unit-ideal quotient. Field-valued points are handled at their actual kernel primes and then by faithful flatness of the field extension, without imposing a closed-point restriction.
\excerpt{PerfectDerivedSupport}{56}{96}

A nonzero field fiber yields some nonzero actual derived cohomology object in an integer degree. The bounded projective model proves vanishing outside its bounds. The induced Hecke action comes from multiplication on the original cochain model, through scalar extension and passage to derived cohomology, and equals the specialized scalar character.

These are formal algebraic conclusions. They neither choose a classical weight nor establish classicality. Revised Claim 9.6 requires the classicality range to contain the model interval; revised 12.4 permits applicability at the selected eigensystem itself. The original unrestricted versions remain partial in their original ledger. The cited external cohomological description of derived support is not needed by this production residue-field support proof.

\section{Artin packets, residual augmentation and finite adequacy}
Hecke packet construction starts with a supplied canonical ring decomposition into actual localizations. The input is \code{ArtinData.Decomposition R}, an \(R\)-algebra equivalence whose coordinates are proved to be the canonical localization maps. It supplies no module direct-sum decomposition and no module localization equivalence; those are constructed subsequently.
\excerpt{ArtinData}{17}{41}

From this input, Lean constructs orthogonal idempotents, their packet submodules, the direct sum, and the packet/localized-module equivalence. It verifies the map \(v\mapsto v/1\). The selected local factor's nilpotent maximal ideal and an invertible generator on every other factor identify the packet with the simultaneous generalized eigenspace for one positive exponent. This does not assume semisimplicity. In characteristic zero, the projector trace is the dimension of its image cast into the field, and nonzero trace detects occupancy. Positive-characteristic trace cancellation remains a real boundary.

Auxiliary group-algebra transfer is constructed on the actual group basis and commutes with actual augmentation. Coinvariant recovery proves a quotient by all balancing relations and identifies it with the ordinary tensor product when the appropriate algebra is commutative. The right-module treatment handles arbitrary groups, including nonabelian two-groups; a supplied free basis descends to a coefficient basis with the same index set, including rank zero.

For a product of cyclic \(p\)-groups in characteristic \(p\), the full multivariable truncated-polynomial algebra equivalence is constructed and its augmentation generators identified. The no-nontrivial-idempotent result covers every finite \(p\)-group, including nonabelian ones, by an orbit-span/fixed-vector argument over the prime field. The dyadic code proves that rational sign projectors are not integral and constructs \(\mathbb F_2[C_2]\cong\mathbb F_2[X]/(X^2)\). None of this supplies dyadic arithmetic image, local geometry or classicality hypotheses.

The \(S_3\) example over \(\mathbb F_7\) is a Lean certificate for one finite representation. The modules construct the standard representation from the permutation module, prove absolute irreducibility, identify its trace-zero adjoint with sign plus standard, and connect averaging to actual first group cohomology. Canonical generalized-eigenspace projectors detect every nonzero adjoint constituent. Independent finite computations can check the printed matrices and trace pairings; their scope remains separate from the formal certificate and from arithmetic adequacy applications. The adequacy convention remains a cited input.

\section{Occupancy, stable cores and adic completion of classes}
In one finite-dimensional carrier, nested idempotent ranges stabilize; nonzero ranges at every level yield a nonzero intersection. In changing carriers, the inverse limit is controlled by the actual intersections of transition images, not merely by nonzero individual levels. \module{StableCores} uses a \code{LinearSystem} with actual linear maps and identity/composition laws. It proves that transition images are antitone, that finite-dimensionality makes each stable core an eventual image, and that the core transition is surjective. Recursive lifting then produces a compatible nonzero vector above a chosen nonzero core coordinate.
\excerpt{StableCores}{6}{38}

Finite companion forests give branches when every finite level is nonempty. An empty level obstructs complete realizations whose actual finite reductions belong to those declared levels; it does not rule out representations excluded by an arbitrary admissibility restriction. The production \code{empty_level_no_realization} explicitly takes reduction maps from the supplied realization type into every forest level. Forests do not alone produce continuous integral representations with every arithmetic condition. Revised Claim 10.5 makes the missing effectivity inputs explicit: complete separated coefficient rings for finitely generated ideals, finite residue rings, finite occupied classes of actual continuous quotient representations, reduction compatibility, continuous coefficient maps, integral Frobenius polynomials with matching scalar changes, and declared polynomial ideal conditions valid in every admissible frame.

Finite frames turn finite admissible conjugacy classes into finite sets of actual representatives. Compactness selects strictly compatible representatives without assuming transition surjectivity or lifting a prescribed conjugating matrix. Entrywise reconstruction, including inverse entries, constructs the actual matrix-group inverse limit. Continuity is derived from finite quotient continuity and the adic topology. Separatedness upgrades compatible polynomial and subgroup-triviality congruences to exact integral equalities. Evaluation then kills each declared polynomial ideal and factors through its actual quotient.
\excerpt{ConditionalCompletion}{8}{54}

The conclusion retains a residual conjugacy class, not a prescribed residual matrix frame; uniqueness up to integral conjugacy is not asserted. No arbitrary Hodge or Weil--Deligne condition is identified with a polynomial ideal without additional arithmetic work. Integral Frobenius lifts and occupancy at every precision are hypotheses, not outcomes of abstract compactness. A missing coefficient-place member is not supplied by propagation inside an existing compatible family.

\section{Conditional realization and the five separate gates}
The formal realization theorem compares actual matrix representations via an invertible linear intertwiner. It assumes a Hausdorff characteristic-zero topological field with continuous ring operations, continuous compared representations, a supplied semisimplification preserving characteristic polynomials, arithmetic density of the declared Frobenius conjugacy saturation, and a general algebraic Brauer--Nesbitt identification theorem over that coefficient field.

\module{DenseFrobenius} proves continuity of characteristic-polynomial coefficients and extends equality from the dense Frobenius conjugacy set to all group elements. The supplied Brauer--Nesbitt input then identifies the semisimple representations. Neither density nor the general Brauer--Nesbitt theorem is a newly introduced global Lean axiom: each is an explicit ordinary theorem argument, and the production result is conditional on it.
\excerpt{GaloisRealization}{13}{53}

For a supplied universal representation, \code{specialize} is literally coefficient change into the chosen field. Its continuity is explicit. The two representations must have the same specialized Hecke polynomials on the same dense good set. The ordinary completion theorem has the universal specialization polynomial identity as a separate explicit premise; its classicality input supplies the attached representation's matching polynomial identity. The final traditional definition 12.1 spells out these inputs; a nonzero fiber and an attached representation alone would not identify the universal specialization.
\excerpt{GaloisRealization}{73}{97}
\excerpt{OrdinaryCompletion}{171}{213}

The five-gate modules derive finite-level support from the actual patched quotient/control data. A general comparison can use an actual inclusion of support subsets under a ring map, with ordinary or derived support as appropriate. The compatible Hecke point \(y\) satisfying \(y\circ f=x\) is separately supplied. Support preservation alone does not provide such a point over the prescribed coefficient field: \(\mathbb Q\to\mathbb Q(i)\) preserves full support but admits no extension of the identity point to \(\mathbb Q\).

The chosen nonzero tensor vector or cohomology vector carries all prescribed Hecke eigenvalues. Classicality is applied to this actual eigensystem, and supplies an occurring automorphic object plus its attached representation. The representation comparison follows under the common-polynomial, continuity, density and semisimplification hypotheses. Propagation uses two already existing compatible families, the family attached to the same automorphic object, distinguished member identities, and their common algebraic spine.

\begin{longtable}{p{.20\textwidth}p{.69\textwidth}}
\toprule Gate & Input and formal consequence\\\midrule\endhead
Residual & Arithmetic residual occupancy and representability supply the deformation problem, its ring and universal representation.\\
Patching & An applicable arithmetic theorem supplies patched carriers, control maps and full patched support. Lean descends support to the actual finite-level carrier.\\
Comparison & A support-preserving comparison and compatible Hecke point are supplied. Lean transports support and constructs a nonzero specialized eigenvector.\\
Fiber & Ordinary tensor or derived cohomology is constructed. The actual Hecke action becomes the specialized scalar character.\\
Realization & Applicability of classicality to the selected eigensystem, occurrence, attachment, common polynomials, density, semisimplification and Brauer--Nesbitt are supplied. Lean deduces the intertwiner.\\
\bottomrule
\end{longtable}

\section{External arithmetic inputs and editorial boundaries}
The formal code begins after the relevant arithmetic objects have been supplied. Retained external inputs include representability and local quotient existence; universal arithmetic deformation representations; residual automorphic occupancy; applicability of a named patching package; classicality at the selected eigensystem; an occurring automorphic object and continuous semisimple attached representation; matching Frobenius normalizations; already existing attached compatible families; arithmetic Frobenius density; construction of semisimplification; and general original-field Brauer--Nesbitt. The canonical Artin localization decomposition and the adequacy convention are also supplied inputs.

An application must state a finite ramification set and coherent choices of Frobenius convention and polynomial normalization. The literal abstract Lean group is \(\Gamma\); its arithmetic realization as \(G_{F,S}\), and density of the corresponding permitted conjugacy classes, are provided by the application. Finite-quotient Chebotarev followed by continuity is the relevant density mechanism for an arbitrary declared Hausdorff coefficient topology. Claims of historical novelty or citation support are not certified by the kernel.

The final traditional manuscript clarifies the statement context without changing production Lean: Claim 3.3 exposes the local structural map and the complete Hausdorff local coefficient-algebra point category; the proof of 6.4 treats \(r>n\) explicitly; and Definition 12.1 exposes the same Frobenius polynomials and all realization hypotheses needed by 9.4. Compatible carrier actions and the declared-level scope of companion obstruction are also explicit, as are the finite ramification set and Frobenius convention. These are final editorial clarifications, tracked separately from the eight inherited amendments. They are not fresh kernel certification of a newly edited proposition.

The strongest global conclusion is conditional automorphic realization inside an explicitly declared supported and realized class, followed by propagation inside already existing compatible systems. Universal residual companions, arbitrary compatible-family existence, unrestricted reciprocity and an arithmetic theorem at coefficient prime two are not conclusions of these sources.

\section{Replay prerequisites, pins and audit procedure}
The production proof replay is Linux/WSL oriented. Native Windows behavior is not claimed. Obtain the separate portable runtime whose root contains \code{toolchain/bin/lean} and \code{workspace/}. It must include the matching cached dependency objects; the guide does not redistribute Lean or the Mathlib dependency cache. Python 3, Bash, Git and ordinary POSIX process tools are needed by the inherited verifier.

\begin{tabular}{p{.23\textwidth}p{.67\textwidth}}
\toprule Component & Required pin\\\midrule
Lean version & \code{4.34.0-rc2}, x86\_64 Linux Release\\
Lean commit & \code{6a10ac8c22beadecabdbb0919c2b50214762f91d}\\
Compiler SHA-256 & \code{e8baaa71855a616dc351028f3ad2200051b0671f423a1696a100e809302d5550}\\
Mathlib commit & \code{aeb72f1e9981b6d35b426b16dca7759be3e7ba06}\\
Options & One thread; warnings as errors; per-process timeout 180 seconds by default\\
\bottomrule
\end{tabular}

Before replay, verify the compiler's exact hash/commit and every dependency HEAD against the runtime's dependency-pin record, not just the version substring and Mathlib HEAD. The inherited verifier checks a version substring, Mathlib HEAD/manifest, production preservation hashes, exact source manifest, object hashes during resume and all axiom targets. It sets \code{MALLOC_ARENA_MAX=2} and \code{-j1}. It does not itself impose an address-space resource limit. To reproduce the reported four-GiB limit on Linux/WSL, enforce it in the caller:
\begin{Verbatim}[fontsize=\small,breaklines=true,breakanywhere=true,frame=lines]
cd formal_release
sha256sum /path/to/Komorebi_Portable_Lean_4.34.0-rc2/toolchain/bin/lean
ulimit -v 4194304
bash Run-Reverified.sh /path/to/Komorebi_Portable_Lean_4.34.0-rc2
\end{Verbatim}
The address-space limit is in KiB in Bash on Linux. Run in a writable extracted copy: replay cleans and rewrites \code{build/}, \code{logs/}, \code{verification.json}, the current axiom consumer and the supplemental consumer report. Preserve the inherited records separately before executing. The Bash wrapper uses bounded four-module batches and resumes only while status is \code{RUNNING}. It rejects a final result other than \code{PASS}.

The cumulative verifier builds every registered production module from source, then imports each exact object and runs \texttt{\#print axioms} for every public named declaration selected by its inventory. Only \code{propext}, \code{Classical.choice} and \code{Quot.sound} are allowed in the transitive audit. It finally tests three invalid inputs: a false proof, a placeholder and an incorrect empty-family normalization. All three must be rejected with the expected diagnostics; a timeout is not a successful negative control.

\code{check_final_claims.py} then verifies source and object hashes against the successful cumulative report and compiles the three direct consumer theorems. Core consumers are \code{PaperXIIFinalConsumer.exact_completed_kernel}, \code{claim31} and \code{claim33}. The native v0.17.1 archive additionally retains a distinct three-theorem \code{PaperXII.ManuscriptContractReplay} consumer and its \code{CONTRACT_REPLAY.json}. Both are supplementary interface evidence, not six additional production theorems and not new numbered claims. The native and core production source bytes agree.

The core inherited cumulative run is recorded at 03:48:31--04:00:23 UTC on September 12, 2026. The native cumulative run is separately recorded at 04:04:33--04:15:36 UTC, with native contracts at 04:16:12--04:16:15 UTC. Matching sources, archived objects and log/report hashes corroborate provenance; they do not transform those historical runs into a new October run.

\section{Package contents, Markdown files and rebuilding this document}
Complete production source is in \code{formal_release/lean/}; it includes all reused dependencies. The directory also supplies the registered manifest, replay scripts, consumers, invalid controls, preserved compiler/axiom logs, inherited verification records, preservation ledgers and required predecessor inventory records. The original raw production sources are unchanged. The native contract consumers and their separate provenance are in \code{native_contract_evidence/}.

The \code{data/} folder supplies JSON and CSV forms of the complete source catalog, declaration census, direct import audit and revised 49-claim crosswalk. Every catalog record includes its full source hash. \code{SOURCE_SHA256SUMS.txt} hashes all production sources. \code{excerpts/} contains 24 full exact core modules, their documented typography copies, and the excerpt index with source hashes. These exact files allow every excerpt's ambient hypotheses and proof text to be inspected without relying on the PDF.

\textbf{Markdown (.md) files are included.} \code{README.md} explains the guide package, evidence scope and build commands; \code{SCOPE_AND_AUDIT.md} explains counts, assumptions and final self-review. Preserved release Markdown notes describe the September replay, final two claims, historical failures and mathematical self-review. Markdown is plain text and can be read without TeX or Lean.

Build this guide from an extracted package using Python 3.10 or later, TeX Live with XeLaTeX, \code{fontspec}, \code{unicode-math}, \code{fvextra}, \code{microtype}, \code{booktabs}, \code{longtable}, \code{fancyhdr} and \code{hyperref}, and Poppler's \code{pdfinfo}, \code{pdftotext}, \code{pdffonts}. Fonts are supplied under relative \code{fonts/} paths with license notices. No machine-specific font path is required.
\begin{Verbatim}[fontsize=\small,frame=lines]
python3 build.py --output-dir /absolute/path/to/fresh-output
\end{Verbatim}
The builder stages its relative inputs in the fresh output directory, fixes \code{SOURCE_DATE_EPOCH} and the document date, suppresses optional timestamp/trailer metadata, and runs three TeX passes. It checks missing glyphs, unresolved references, overfull boxes, required text anchors and embedded fonts. Its report records PDF, extracted-text and source hashes plus page count. Independent fresh-directory rebuilds must agree byte-for-byte on the PDF, text and pages under the documented same tool versions. PDF determinism is distinct from proof replay determinism.

\appendix
\section{Full source catalog}
The complete 96-module production catalog follows in registered dependency order. ``Thm.'' counts each module's public theorem declarations, including reused declarations where present. SHA prefixes aid visual comparison; the exact 64-character hashes and imports are supplied in \code{data/source_catalog.json} and \code{data/source_catalog.csv}. Source paths are determined by replacing dots in the module name with slashes and adding \code{.lean}, relative to \code{formal_release/lean/}.
{\small\begin{longtable}{p{.37\textwidth}r r p{.35\textwidth}}
\toprule Module & Lines & Thm. & Source SHA-256 (prefix)\\
\midrule\endhead
\code{PaperX.Reconstruction} & 189 & 12 & \texttt{d8de70c175635ea47d8eaed4}\\
\code{PaperXII.ArtinData} & 97 & 9 & \texttt{953b639b0309f883d326d0f8}\\
\code{PaperXII.AuxiliaryGroupAlgebra} & 72 & 7 & \texttt{ef811535231880a3b7aeb984}\\
\code{PaperXII.CompletionQuotient} & 181 & 19 & \texttt{c7b60b3b5a59cf73ce1891b7}\\
\code{PaperXII.ComplexBaseChange} & 86 & 3 & \texttt{6d91bbc8a4849ac4a5bae4f9}\\
\code{PaperXII.DenseFrobenius} & 113 & 8 & \texttt{941465b7bae2df984810d17a}\\
\code{PaperXII.DeterminantalIdeals} & 131 & 11 & \texttt{1091c57b8583977120bbe9d2}\\
\code{PaperXII.DyadicProjectors} & 195 & 17 & \texttt{a2d01c1ddf2152e2b6bd0269}\\
\code{PaperXII.FiniteCarrier} & 67 & 7 & \texttt{ccc9770f4c6cc72b83a94f0a}\\
\code{PaperXII.FiniteGroupCohomology} & 107 & 9 & \texttt{abab8fb127de1bbc804296e2}\\
\code{PaperXII.FiniteTensorQuotient} & 172 & 15 & \texttt{1ede310ac1f69aeeab66b851}\\
\code{PaperXII.LocalIdempotents} & 18 & 1 & \texttt{e34c194a98e51680b23542d6}\\
\code{PaperXII.LocalQuotients} & 126 & 11 & \texttt{738ec11826cea7e7aabccb35}\\
\code{PaperXII.ModuleSupport} & 63 & 8 & \texttt{b724c27d7ea2d3f2a5941ffe}\\
\code{PaperXII.PacketTrace} & 74 & 8 & \texttt{ad4c7192632b5075026bb73c}\\
\code{PaperXII.RepresentationData} & 86 & 6 & \texttt{6f938f551619e48f1f5a94ae}\\
\code{PaperXII.ReuseCompactness} & 175 & 12 & \texttt{0bd9169cb2b0efc4938e70d7}\\
\code{PaperXII.ReuseFiniteSeparation} & 97 & 6 & \texttt{1975b5fb629823eb7f80c17c}\\
\code{PaperXII.S3Standard} & 221 & 22 & \texttt{c483d46bda60362ef0650d78}\\
\code{PaperX.RepresentationLimits} & 193 & 15 & \texttt{81fb37746e7470557deced36}\\
\code{PaperXII.AdicModel} & 73 & 6 & \texttt{61c8a69111fc54671eb20b99}\\
\code{PaperXII.ArtinModules} & 206 & 15 & \texttt{78d81338b8f8f57638051245}\\
\code{PaperXII.AugmentationCoinvariants} & 120 & 8 & \texttt{9bbc93c64a5beccbe8fbc440}\\
\code{PaperXII.CompanionForests} & 84 & 10 & \texttt{db2d885e73099916d65bcaf0}\\
\code{PaperXII.ComplexFieldFibers} & 108 & 9 & \texttt{cf9b3255fcadbb9d1807f957}\\
\code{PaperXII.ComponentSupport} & 21 & 1 & \texttt{d6a2a42f8834ba233d112609}\\
\code{PaperXII.DeterminantalStabilization} & 130 & 10 & \texttt{7ee60aa621c1a1b775918634}\\
\code{PaperXII.ExistingCompatibleSystems} & 81 & 3 & \texttt{ed9748c8de4dc8693af79eb9}\\
\code{PaperXII.FieldFibers} & 73 & 6 & \texttt{36d276ab57e7cc632e15571f}\\
\code{PaperXII.FrobeniusSeparation} & 76 & 7 & \texttt{009a953c292806afb770c631}\\
\code{PaperXII.GaloisRealization} & 98 & 5 & \texttt{19e69eedbf482775e718e525}\\
\code{PaperXII.GeneralizedPackets} & 137 & 7 & \texttt{4a67970524a245b0414d0b7c}\\
\code{PaperXII.NonabelianResidual} & 142 & 8 & \texttt{923412037be8c47c1017e742}\\
\code{PaperXII.PerfectModels} & 122 & 2 & \texttt{fa7dd21c8776d96b8ab917cc}\\
\code{PaperXII.PowerSeriesComponent} & 246 & 38 & \texttt{cd0599c3aee9d2397d505928}\\
\code{PaperXII.QuotientSupport} & 60 & 5 & \texttt{733d40c9b25bbee513955de0}\\
\code{PaperXII.ResidualAugmentation} & 78 & 6 & \texttt{8ceac287f565bf81c0267dae}\\
\code{PaperXII.S3Adjoint} & 240 & 16 & \texttt{c5cee6cdd374926c9897e0a6}\\
\code{PaperXII.S3Detection} & 316 & 26 & \texttt{365b36245a4c4aa11c4522a4}\\
\code{PaperXII.S3Spectrum} & 86 & 8 & \texttt{ca1cdfc3a4797bc71ee146fa}\\
\code{PaperXII.StableCores} & 92 & 6 & \texttt{e266cbe63a32d7c4d3b719a3}\\
\code{PaperX.RepresentationClasses} & 132 & 16 & \texttt{a35e3a56c160992d4883ced0}\\
\code{PaperX.RepresentationTopology} & 124 & 8 & \texttt{9952178e10c02b66f27949cd}\\
\code{PaperXII.AdicQuotientTopology} & 105 & 10 & \texttt{c3ffb8523760a84befb64bb0}\\
\code{PaperXII.BalancedAugmentation} & 128 & 9 & \texttt{dd128eb806a101cebe777f9a}\\
\code{PaperXII.BaseChangeSupport} & 47 & 4 & \texttt{0da87b182a1e79db32ef36c5}\\
\code{PaperXII.ComplexSupport} & 103 & 7 & \texttt{186a3d1126dfa9296836b763}\\
\code{PaperXII.DeterminantalBlockStabilization} & 83 & 8 & \texttt{f6657438b1c2b5fd31012d09}\\
\code{PaperXII.DyadicProjectorsResidual} & 124 & 13 & \texttt{757fa5f1c294698c32648da8}\\
\code{PaperXII.FiniteAdequacy} & 109 & 6 & \texttt{aa7fba126405fa95fcb12944}\\
\code{PaperXII.InverseLimitRepresentations} & 169 & 9 & \texttt{ba425975d3baba88d65b5316}\\
\code{PaperXII.OrdinaryPatching} & 98 & 5 & \texttt{5724b6d1587cb4c657d25380}\\
\code{PaperXII.SupportTransport} & 59 & 4 & \texttt{717228feacde958e8dcd7638}\\
\code{PaperXII.TruncatedGroupAlgebra} & 297 & 25 & \texttt{676626c8a833d93839711162}\\
\code{PaperXII.VectorInverseLimit} & 122 & 10 & \texttt{ae0a5eff607a14a7e93fd523}\\
\code{PaperX.ClassInvariants} & 103 & 12 & \texttt{0506bcb396c1d8cfce4dfbed}\\
\code{PaperXII.AdicCompletionComparison} & 210 & 14 & \texttt{2dd1070d880040a9990f96b3}\\
\code{PaperXII.AdicMatrixGroups} & 163 & 8 & \texttt{21d738f44b4c35e4778528e0}\\
\code{PaperXII.CoinvariantTensorComparison} & 108 & 4 & \texttt{e920619c2a26f8b003123021}\\
\code{PaperXII.CompletedFiniteTensor} & 134 & 7 & \texttt{9c4e04f241b588d50db1d905}\\
\code{PaperXII.CompletedGlobalBaseChange} & 263 & 18 & \texttt{11fc86e509581906275b565e}\\
\code{PaperXII.CorrectedOccupancy} & 24 & 1 & \texttt{efbbb72e4f864a073569846e}\\
\code{PaperXII.HeckeFibers} & 90 & 7 & \texttt{3a0e3363a03013bd87c2f504}\\
\code{PaperXII.PerfectDerivedSupport} & 119 & 9 & \texttt{f0e29718343af1845da6b39b}\\
\code{PaperXII.PresentationComparison} & 122 & 4 & \texttt{f64299848f002c18396ce9bd}\\
\code{PaperXII.RightAugmentation} & 91 & 3 & \texttt{5fdf4abbf1666a8e03ab4cd9}\\
\code{PaperXII.AdicCompanions} & 134 & 7 & \texttt{b1f7e850bd067860825671d1}\\
\code{PaperXII.Coinvariants} & 47 & 2 & \texttt{7500fa7b2277651c91f93365}\\
\code{PaperXII.CompletedBaseChangeTransport} & 113 & 7 & \texttt{332e73f093c5fefc16ad077d}\\
\code{PaperXII.DerivedHeckeAction} & 116 & 7 & \texttt{b1cb976fcc30917465dca404}\\
\code{PaperXII.DerivedPatching} & 92 & 8 & \texttt{853d09528413ec978e2ea106}\\
\code{PaperXII.FiniteTensorLocality} & 149 & 11 & \texttt{be33c445a4329462b7039912}\\
\code{PaperXII.OrdinaryCompletion} & 213 & 9 & \texttt{07087a7da30fd656baed57b3}\\
\code{PaperXII.PerfectCohomology} & 108 & 10 & \texttt{2c560e3c14cc3082ffeb9ef4}\\
\code{PaperXII.PresentationIndependence} & 63 & 2 & \texttt{88a4c15e4ba774f041eeb364}\\
\code{PaperXII.AdicSpines} & 112 & 6 & \texttt{f90d8f75e4106eacda34b776}\\
\code{PaperXII.CompletedGlobalQuotient} & 173 & 17 & \texttt{71eede618e5c6aa481d5ec6b}\\
\code{PaperXII.CompletedTensorUniversal} & 99 & 7 & \texttt{92a58d870df81ff11bdc3f97}\\
\code{PaperXII.DerivedRealization} & 177 & 7 & \texttt{0316cfac39230ec302b3ae8b}\\
\code{PaperXII.HeckeCarriers} & 237 & 16 & \texttt{dd3bdac5d81c73a6d28bb62b}\\
\code{PaperXII.IntrinsicFitting} & 69 & 4 & \texttt{33c122a7c407120e6bda7a36}\\
\code{PaperXII.CompletedLocalConditions} & 100 & 3 & \texttt{7f7f43315592150c1f9b0e48}\\
\code{PaperXII.ContinuousClasses} & 112 & 10 & \texttt{45111ecf0dc7bd82f36a468d}\\
\code{PaperXII.FittingBaseChange} & 60 & 3 & \texttt{5d693d829097c7810a8c3ef6}\\
\code{PaperXII.PatchingCarriers} & 109 & 5 & \texttt{3c8506f8babeddb508dbeb92}\\
\code{PaperXII.ClassCompanions} & 197 & 15 & \texttt{9b177812cf70d0fc68998853}\\
\code{PaperXII.CompletedPaperCapstone} & 69 & 3 & \texttt{2489bf5ea48d25a39a90ea28}\\
\code{PaperXII.FittingTransport} & 60 & 3 & \texttt{23d49ac97ae28abf066ebd10}\\
\code{PaperXII.FiveGateCompletion} & 168 & 5 & \texttt{f107470686ed79b1cf37952c}\\
\code{PaperXII.ClassSpines} & 116 & 4 & \texttt{00aed3a592da86f63b02a21a}\\
\code{PaperXII.CompletedPointCategory} & 101 & 5 & \texttt{1e1fe94ece8611b4b6496b82}\\
\code{PaperXII.FittingBoundaryChecks} & 66 & 6 & \texttt{8516575908e9f9351e9f4666}\\
\code{PaperXII.AdicLocalConditions} & 116 & 8 & \texttt{09cce645d5f4f86c36db909a}\\
\code{PaperXII.CompletionBoundaryChecks} & 45 & 4 & \texttt{ed6c4701f20500afe72ba1c8}\\
\code{PaperXII.FittingStrata} & 64 & 5 & \texttt{a5313b7c1341dbf6bd573598}\\
\code{PaperXII.ConditionalCompletion} & 54 & 1 & \texttt{cc5234bb8a3172352e4b8e70}\\
\bottomrule\end{longtable}
}

\section{All 49 revised internal claims}
The crosswalk lists every revised internal claim, its manuscript label, amendment status, original status and number of mapped targets. Every target name and source module is supplied in \code{data/claim_crosswalk.json}; the condensed module crosswalk is \code{data/claim_crosswalk.csv}. ``Amended'' means complete for the revised statement while the original statement retains its own status. Ordinary editorial clarifications described in Section 13 are separately tracked.
{\small\begin{longtable}{l p{.34\textwidth} l l r}
\toprule Claim & Manuscript label & Revised & Original & Targets\\
\midrule\endhead
3.1 & \code{thm:tensor_quotient} & complete & complete & 17\\
3.2 & \code{cor:simultaneous_local} & complete & complete & 3\\
3.3 & \code{prop:global_local_basechange} & complete & complete & 13\\
3.4 & \code{prop:local_idempotent} & complete & complete & 1\\
4.2 & \code{thm:s3_adequacy} & complete & complete & 20\\
5.1 & \code{prop:group_algebra_transfer} & complete & complete & 5\\
5.2 & \code{thm:residual_augmentation_algebra} & complete & complete & 7\\
5.3 & \code{cor:no_residual_character_splitting} & complete & complete & 7\\
5.4 & \code{thm:coinvariant_recovery} & complete & complete & 6\\
6.2 & \code{thm:finite_module_support} & complete & complete & 1\\
6.3 & \code{thm:fitting_support} & complete & complete & 2\\
6.4 & \code{thm:higher_fitting_strata} & complete & complete & 2\\
6.5 & \code{prop:fitting_base_change} & complete & complete & 2\\
6.7 & \code{thm:near_faithful_full_support} & complete & complete & 4\\
6.8 & \code{cor:every_component_support} & complete & complete & 2\\
6.9 & \code{thm:quotient_support} & complete & complete & 3\\
6.11 & \code{thm:ordinary_support_descent} & complete & complete & 4\\
6.12 & \code{cor:componentwise_occupancy} & amended & unchecked & 1\\
6.16 & \code{thm:derived_quotient_support} & complete & complete & 20\\
6.18 & \code{thm:derived_support_descent} & complete & complete & 15\\
6.19 & \code{cor:derived_degree_survives} & complete & complete & 7\\
7.1 & \code{thm:rt_support_transport} & complete & complete & 2\\
7.2 & \code{cor:rt_full_support} & complete & complete & 1\\
7.3 & \code{prop:quotient_erases_component} & complete & complete & 10\\
7.4 & \code{thm:field_fiber_criterion} & complete & complete & 2\\
7.5 & \code{cor:supported_hecke_fiber} & complete & complete & 1\\
8.2 & \code{thm:canonical_hecke_packets} & complete & complete & 11\\
8.3 & \code{cor:simultaneous_generalized_eigenspaces} & complete & complete & 6\\
8.4 & \code{cor:rank_trace_detection} & complete & complete & 6\\
8.7 & \code{thm:nested_finite_carrier} & complete & complete & 1\\
8.8 & \code{cor:finite_witness_incompatibility} & complete & complete & 1\\
8.10 & \code{lem:stable_core_surjectivity} & complete & complete & 1\\
8.11 & \code{thm:nonzero_stable_core} & complete & complete & 2\\
8.12 & \code{cor:surjective_inverse_system} & complete & complete & 1\\
9.3 & \code{thm:support_gives_hecke_system} & complete & complete & 2\\
9.4 & \code{thm:pointwise_galois_realization} & amended & partial & 7\\
9.5 & \code{cor:full_support_pointwise_automorphy} & amended & partial & 1\\
9.6 & \code{thm:derived_pointwise_realization} & amended & partial & 21\\
10.3 & \code{thm:companion_branch_criterion} & complete & complete & 2\\
10.4 & \code{cor:empty_companion_obstruction} & complete & complete & 2\\
10.5 & \code{thm:simultaneous_companion_completion} & amended & partial & 1\\
10.7 & \code{thm:one_member_propagation} & complete & complete & 2\\
10.9 & \code{prop:finite_frobenius_separator} & complete & complete & 2\\
11.1 & \code{prop:dyadic_sign_projectors} & complete & complete & 14\\
11.2 & \code{cor:dyadic_coinvariant_recovery} & complete & complete & 5\\
12.2 & \code{thm:ordinary_conditional_completion} & amended & partial & 7\\
12.4 & \code{thm:derived_conditional_completion} & amended & partial & 14\\
12.5 & \code{thm:system_completion} & complete & complete & 1\\
12.6 & \code{thm:five_gate_completion} & amended & partial & 36\\
\bottomrule\end{longtable}
}

\section{A practical reading and inspection sequence}
A reader examining a mathematical claim should first locate its revised number in the crosswalk, then read the exact production declarations, their ambient variables and the relevant proof. Follow internal imports to any structure or construction the theorem consumes. Determine which fields are supplied data and which declarations construct the desired maps, objects or identities. Match the point category, coefficient maps, topologies, completeness, finite-generation and locality assumptions with the intended application. Consult the original ledger when comparing against earlier wording.

For Claims 3.1 and 3.3, the central source chain is the raw tensor quotient, completion/closed quotient comparison, adic quotient topology, completed finite tensor, algebraic adic comparison/locality, global tensor transport, local condition factorization, Noetherian/local capstone and coefficient/local point category. Checking only ordinary tensor quotients or an assumed universal-property interface would miss the constructed completed objects and their actual topologies.

For the arithmetic conclusion, trace the nonzero fiber or cohomology vector through the actual carrier and point. Verify classicality for that selected eigensystem, then check attachment, matching Frobenius polynomials, density, semisimplification and original-field representation comparison. Finally ensure any propagated family already exists and belongs to the same attached automorphic object. This sequence explains exactly what the formalization establishes and what an arithmetic application must still provide.
\end{document}
